\chapter*{Introduction}
Drug discovery\cite{bib32} is an arduous journey with many unknowns throughout the path, with several traps that lead to failure. To perform a large number of trials in an iterative manner, computer simulations\cite{bib33, bib34} are convenient to perform tens of thousands of calculations to provide crucial insights, and possibly, rule out dead ends. 
Structure-based drug discovery (SBDD) is one such approach that aids in designing and optimizing chemical entities by leveraging the tertiary structural information of biological targets. SBDD encompasses various strategies, including virtual screening, where large libraries of compounds are screened for its binding affinity with the target site. This process is also called as “hit identification”, where hits are referred to as molecules that show weak yet measurable binding affinity. Hits generally contain most of the chemical framework capable of showing the desired pharmacological response, however, they need optimisation to improve the binding affinity, and other physicochemical properties essential for developing a clinical candidate. Most of the docking programs have shown great success at the lead modification stage, which involves evaluating the effect of the new substituents on the same central core (also called a parent nucleus)\cite{bib53}.

Another important method, called as \textit{de novo} drug design, involves constructing new molecules from molecular fragments within the binding site of the target.  Molecular docking\cite{bib35} is an important technique in structure-based drug discovery (SBDD) that predicts optimal binding conformation of a small molecule (ligand) within the binding site of a biological target. This process helps identify potential drug candidates by estimating their binding affinity and types of interactions to receptor atoms. The latter information is used to suggest modifications in the hit molecules, such that their binding affinity is improved. Iterative cycles of design, synthesis, and experimental validation refine these compounds to optimize their pharmacological properties, this process is called as “lead optimisation”. 

With a near exponential increase in the tertiary structure information in the last 3 years attributed to several artificial intelligence based methods, particularly Alphafold\cite{bib37} and RoseTTAFold\cite{bib38}, many efforts have been made to incorporate machine learning methods in drug discovery to improve the quality of therapeutic agents and reduce the drug discovery timelines significantly\cite{bib36}. Furthermore, developers are also constantly developing methods that employ machine learning methods to improve the prediction accuracy of molecular docking methods. Incorporating the state-of-the-art deep learning neural networks have shown to improve conformational search algorithms and develop better and more generalised scoring functions. These improvements have helped ameliorate challenges posed by limited structural and small molecule data\cite{bib39, bib40}. 


Molecular docking has been used for more than four decades that serves to attain two primary objectives:
\begin{itemize}
    \item To predict the binding affinity and conformation of small molecules within a receptor site, and 
    \item To identify hits from a large chemical database to search for diverse chemical scaffolds.\cite{bib1}. 
\end{itemize}

Despite classifying these objectives behind molecular docking into separate classes, their boundaries are not clearly demarcated. To attain these objectives, molecular docking programs must comprise components that can search the conformational space within the binding site and a scoring function that scores the conformations such that the biological relevant ones are ranked higher.
